% s-223-p223-m002; original PDF 223, 224
\Aphorism{s-223-p223-m002}{316}
\Sperr{Der Rationalist redet. }Es wäre möglich, dass ein entwickelter Geschmack das Drama, im älteren drastischen Sinn, noch einmal gänzlich hinter sich liesse, aus dem einfachen Grunde, dass er die Menschen dieses Dramas nicht mehr ertragen kann. Wenn wir eine Shakespeare’sche Tragödie hören, oder eine griechische, oder eine von Kleist, so sind wir löblich im evangelischen Sinne; unsere Linke weiss nicht, was die Rechte thut — offen geredet: wir üben ein Mass von Nachsicht, fühlen uns aber nicht als Verzeihende, sondern sogar als Verehrende. Welche Art Mensch lassen wir uns in so einem blutrünstigen Bühnenstück bieten! Intellecte, die ihre viereckige Beschränktheit im Triumph dahertragen wie der Ochse sein Kummet; Temperamente ohne Stil, dumpf, chaotisch, blindwüthig; Charaktere, die sich, wie der Hund in den Knochen, in eine einzige Kategorie des moralischen Lebens verbeissen und davon nicht ablassen; eine Halsstarrigkeit, die jeden Versuch der Aufklärung erstickt; ein Brüllen und Toben der Leidenschaften, dass Wort und Gedanke überhört werden — eine vorsintfluthliche brutale, riesenhaft-tölpische Art Mensch, der wir in der Gesellschaft hundert Schritt aus dem Wege gehen würden. Und warum dulden wir sie auf der Bühne? als Mittel zum Zweck, einen Theaterabend zu ermöglichen. Diese Stücke wären in der ersten Scene zu Ende, wenn die Vernunft einen Augenblick zu Gehör käme, wenn sich die Betheiligten über das einigen könnten, worüber sie sich schliesslich im fünften Act, nach einigen Stunden Verschwörung, Blutfliessen und Todtschlag, einigen müssen. Diese Stücke illustriren alle den Satz, dass der pathetische Mensch dumm ist und über Kieselsteine stolpert, dass dort, wo der vernünftige Mann \Foreign{prae limine} zaudert und ablässt, die Theaterfigur darauf los wüthet und erst durch das geschehene Unheil klug wird, dass auf der Bühne der indirecte Beweis, die \Foreign{deductio ad absurdum,} selbst dort nöthig ist, wo im Leben alles sonnenklar wäre. Aber diese künstliche Vernagelung der Köpfe und Verfinsterung der Herzen ist, wie gesagt, Vorbedingung des Dramas; wenn König Lear einen Augenblick darüber nachdächte, was Cordelia meint, oder Cordelia dem Alten zu Gefallen eine Nüance liebenswürdiger antwortete, so wäre Conflict und Tragödie im Keime erstickt. Freilich hat der Künstler das Recht zur \Foreign{petitio principii,} man muss einen Knoten machen, ehe man nähen kann — aber diese Dramatik, um im Bilde zu bleiben, macht einen falschen Knoten und muss mit sehr lockeren Stichen nähen, damit er nicht zu früh aufgehe!
\par

